\chapter{TCP, UDP and the Diagnostic Toolkit}\label{ch:transport}
\bp{2.4}

\begin{outcomes}
By the end of this chapter you will be able to:
\begin{tight}
  \item \textbf{Explain} what a port number identifies, and recognise the
        well-known ports that appear in access lists and in fault reports.
  \item \textbf{Choose} correctly between TCP and UDP behaviour when reasoning
        about a symptom.
  \item \textbf{Use} ping, extended ping and traceroute deliberately, varying
        source, size and count to test a specific hypothesis.
  \item \textbf{Read} packet capture output well enough to say which direction
        failed.
\end{tight}
\end{outcomes}

\begin{prereq}
\chref{ipv4} for addressing, \chref{ethernet} for what a switch does with a
frame.
\end{prereq}

\begin{keyterms}
port $\bullet$ socket $\bullet$ TCP $\bullet$ UDP $\bullet$ three-way handshake
$\bullet$ acknowledgement $\bullet$ retransmission $\bullet$ well-known port
$\bullet$ ephemeral port $\bullet$ ICMP $\bullet$ echo request $\bullet$ TTL
$\bullet$ extended ping $\bullet$ traceroute
\end{keyterms}

\begin{examnote}[The toolkit is the exam]
Topic \tp{2.4} names the tools by name: ``show commands (including show logs),
ping, extended ping, trace route, and packet capture output''. It is one of the
eight \emph{troubleshoot} topics, and the tools in it are how you will answer the
other seven. Learn them as instruments, not as trivia.
\end{examnote}

\section{Ports, and what they identify}

An IP address identifies a host. A \term{port number} identifies which program on
that host. Together they form a \term{socket}, and a conversation is defined by
four things: source address, source port, destination address, destination port.

\begin{longtable}{@{}C{1.6cm}L{2.6cm}L{2.0cm}L{8.6cm}@{}}
\toprule
\rowcolor{primary!15}
\textbf{Port} & \textbf{Service} & \textbf{Protocol} & \textbf{Why it matters here} \\
\midrule
\endhead
22  & SSH        & TCP & Device management. \chref{aaa} \\
\rowcolor{lightbg}
23  & Telnet     & TCP & Device management, unencrypted. Disable it \\
53  & DNS        & UDP, TCP & Name resolution. \chref{dns} \\
\rowcolor{lightbg}
67, 68 & DHCP    & UDP & Address assignment. \chref{dhcp} \\
69  & TFTP       & UDP & Legacy file transfer. Replaced by SFTP/SCP in v2.0 \\
\rowcolor{lightbg}
80, 443 & HTTP, HTTPS & TCP & The traffic users notice \\
161, 162 & SNMP  & UDP & Monitoring and traps. \chref{management} \\
\rowcolor{lightbg}
514 & Syslog     & UDP & Log delivery. \chref{management} \\
\bottomrule
\end{longtable}

Client programs use a high-numbered \term{ephemeral} port as their source. This
is why a return packet is not simply ``port 80'' --- it is addressed to whatever
ephemeral port the client opened, and why stateless access lists need care
(\chref{acl}).

\begin{conceptbox}[TCP or UDP, and what each failure looks like]
\textbf{TCP} establishes a connection with a three-way handshake --- SYN,
SYN-ACK, ACK --- numbers every byte, acknowledges what arrives and retransmits
what does not. It costs latency and gains reliability.

\textbf{UDP} sends and forgets. No handshake, no acknowledgement, no
retransmission. It is right for anything where a late packet is worse than a lost
one, and for short exchanges where a handshake would double the cost --- which is
why DNS and DHCP use it.

The diagnostic value is in how they fail. A blocked TCP port produces a
\emph{connection timeout or refusal} --- something definite. A blocked UDP port
often produces \emph{nothing at all}, because there was never going to be an
acknowledgement. ``It just hangs'' points at UDP or at a one-way path; ``connection
refused'' means something answered and said no.
\end{conceptbox}

\section{Ping, and how to ping deliberately}

Ping sends an ICMP echo request and waits for a reply. It proves that a path
exists in \emph{both} directions --- a point worth dwelling on, because a
successful ping rules out far more than a failed one confirms.

\begin{verify}{R1}{ping 192.168.20.10}
Type escape sequence to abort.
Sending 5, 100-byte ICMP Echos to 192.168.20.10, timeout is 2 seconds:
!!!!!
Success rate is 100 percent (5/5), round-trip min/avg/max = 1/2/4 ms
\end{verify}

\begin{longtable}{@{}C{1.6cm}L{13.6cm}@{}}
\toprule
\rowcolor{primary!15}
\textbf{Symbol} & \textbf{Meaning} \\
\midrule
\endhead
\cmd{!} & Reply received. The path works in both directions \\
\rowcolor{lightbg}
\cmd{.} & Timed out. Either no path, no return path, or something is dropping
silently \\
\cmd{U} & Destination unreachable --- a router explicitly told you it has no
route. This is far more informative than a timeout \\
\rowcolor{lightbg}
\cmd{Q} & Source quench: congestion \\
\cmd{M} & Could not fragment. An MTU problem \\
\bottomrule
\end{longtable}

\begin{conceptbox}[The extended ping is the one that finds faults]
A plain ping from a router sources the packet from the exit interface. That is
often not the address you care about. \term{Extended ping} lets you set the
source, and this single option answers the commonest routing question: ``does the
far end have a route \emph{back} to this subnet?''

Run \cmd{ping} with no arguments to enter the extended dialogue.
\end{conceptbox}

\begin{verify}{R1}{ping}
Protocol [ip]:
Target IP address: 192.168.30.10
Repeat count [5]: 100
Datagram size [100]: 1500
Timeout in seconds [2]:
Extended commands [n]: y
Source address or interface: 192.168.20.1
Type of service [0]:
Set DF bit in IP header? [no]: yes
Sweep range of sizes [n]:
Sending 100, 1500-byte ICMP Echos to 192.168.30.10, timeout is 2 seconds:
Packet sent with a source address of 192.168.20.1
!!!!!!!!!!!!!!!!!!!!!!!!!!!!!!!!!!!!!!!!!!!!!!!!!!
Success rate is 100 percent (100/100), round-trip min/avg/max = 1/3/12 ms
\end{verify}

\begin{tight}
  \item \textbf{Source address} --- test the return path for a specific subnet.
  \item \textbf{Repeat count} --- 100 pings exposes intermittent loss that five
        will miss.
  \item \textbf{Datagram size} with \textbf{DF bit set} --- finds MTU problems.
        If 1500 fails and 1400 succeeds, something on the path cannot carry a
        full-size frame.
\end{tight}

\section{Traceroute}

Traceroute sends packets with deliberately small TTL values. The first router
decrements TTL to zero, discards the packet and reports back; the second does the
same for TTL 2, and so on. The result is a list of hops.

\begin{verify}{R1}{traceroute 192.168.30.10}
Type escape sequence to abort.
Tracing the route to 192.168.30.10

  1 10.0.12.2 4 msec 1 msec 1 msec
  2 10.0.23.3 8 msec 5 msec 4 msec
  3  *  *  *
  4  *  *  *
\end{verify}

\begin{alertbox}[Where a traceroute stops is not always where the fault is]
Asterisks mean no reply from that hop. That can mean the fault is there --- or
that the hop is configured not to reply, or that the \emph{return} path from that
hop is broken while the forward path is fine. A traceroute that dies at hop 3
narrows the search to hop 3 onwards; it does not convict hop 3. Confirm with an
extended ping sourced from an address the far end should know.
\end{alertbox}

\section{Reading a capture}

\begin{verify}{SW1}{show capture BUF}
   1  0.000000  192.168.20.10 -> 192.168.30.10  ICMP Echo request  id=1 seq=1
   2  2.001113  192.168.20.10 -> 192.168.30.10  ICMP Echo request  id=1 seq=2
   3  4.002240  192.168.20.10 -> 192.168.30.10  ICMP Echo request  id=1 seq=3
\end{verify}

Three requests, no replies. That is a one-way path: the packets are leaving and
arriving, but nothing is coming back. The fault is on the return path or at the
destination host, and it is emphatically \emph{not} the local switch --- which is
exactly the conclusion the capture is there to let you draw.

\begin{pitfall}
\begin{tight}
  \item \textbf{Treating a failed ping as proof of a broken forward path.} Ping
        needs both directions. Half the time the forward path is fine.
  \item \textbf{Pinging from the wrong source.} A router pinging from its exit
        interface tests a path the users never take.
  \item \textbf{Concluding from five packets.} Intermittent loss needs a hundred.
  \item \textbf{Forgetting that ICMP may be filtered.} Plenty of hosts and
        firewalls drop echo requests by policy. ``It does not ping'' is not the
        same as ``it is down''.
\end{tight}
\end{pitfall}

\begin{breakfix}[Web server reachable by IP, refuses connections by name and port.]
Users cannot load an internal site. From a workstation, \cmd{ping 10.20.5.40}
succeeds with 0\% loss. Browsing to the address times out.

\textbf{Diagnosis.} Ping proves layer 3 connectivity in both directions, so
addressing, masks, routing and the physical path are all sound. The failure is
above layer 3 --- either the service is not listening, or something is filtering
TCP 443 specifically. A timeout rather than a refusal points at a filter: a
closed port on a live host normally answers with a reset.

\textbf{Fix.} Test the port rather than the host. From a router on the path, an
extended ping cannot help --- it is ICMP only --- so use \cmd{telnet 10.20.5.40
443}: a blank screen means the port is open, a refusal means closed, and a hang
means filtered. Then walk the access lists on the path (\chref{acl}). Do not
touch addressing; the ping already cleared it.
\end{breakfix}

\begin{lab}{Diagnose with intent}{Packet Tracer}
\textbf{Topology.} Two routers back to back, one LAN each, one PC per LAN.

\begin{tasks}
  \item Address everything and confirm end-to-end ping between the PCs.
  \item From \cmd{R1}, run a plain ping to the far PC, then an extended ping
        sourced from the LAN interface. Explain why the second is the more
        useful test.
  \item Remove the return route on \cmd{R2}. Predict what a ping from PC-A will
        show, and what a traceroute will show, \emph{before} testing.
  \item Restore the route. Now set the far LAN interface MTU to 1400, and use an
        extended ping with size 1500 and the DF bit set. Record the result.
  \item Run a hundred-packet ping and record whether any loss appears.
  \item \textbf{Extension.} Have a partner break one thing --- an address, a
        mask, a route or an interface --- and identify which, using only ping,
        extended ping and traceroute, in under five minutes.
\end{tasks}
\end{lab}

\begin{checkpoint}
\begin{tightnum}
  \item A ping returns \cmd{U.U.U}. What has happened, and why is that more
        useful than five dots?
  \item Why does an extended ping with a specified source often succeed when a
        plain ping from the same router fails?
  \item A traceroute shows three hops then asterisks. List three different
        explanations.
\end{tightnum}
\end{checkpoint}

\begin{chaptersummary}
\begin{tight}
  \item A port identifies the program; address plus port is a socket.
  \item TCP handshakes, acknowledges and retransmits; UDP does none of it. Their
        failure modes differ, and the difference is diagnostic.
  \item Ping tests both directions at once. A success rules out a great deal; a
        failure convicts nothing on its own.
  \item Extended ping sets the source, size and DF bit --- it is the tool that
        finds return-path and MTU faults.
  \item Traceroute narrows the search; it does not name the culprit.
  \item A capture showing requests and no replies is a one-way path.
\end{tight}
\end{chaptersummary}

\begin{reviewq}
  \item \textbf{(MCQ)} Which pair of ping output characters indicates that a
        router explicitly reported it had no route?
        (a)~\cmd{.....} (b)~\cmd{UUUUU} (c)~\cmd{!!!!!} (d)~\cmd{MMMMM}
  \item \textbf{(MCQ)} Which protocol would you expect a DNS query to use by
        default? (a)~TCP 53 (b)~UDP 53 (c)~UDP 67 (d)~TCP 22
  \item \textbf{(Short)} Explain why setting the source address on a ping is the
        single most useful extended option when troubleshooting routing.
  \item \textbf{(Short)} A service ``hangs'' rather than being refused. What does
        that suggest, and what would you check?
  \item \textbf{(Applied)} A user reports a file transfer that succeeds for small
        files and fails for large ones. Describe, with the exact ping options you
        would use, how you would confirm or rule out an MTU problem.
\end{reviewq}
